% Glossary Appendix
% Auto-generated from glossary.yaml

\chapter{Glossary}

This glossary defines all technical terms used throughout the North Star Document.

\section*{Foundation Systems}

\textbf{atom:} Smallest persistent memory unit (content + tags + metadata + provenance) created by \cmc.

\textbf{\cmc\ (Context Memory Core):} Immutable atom storage system providing durable memory with provenance tracking.

\textbf{\hhni\ (Hierarchical Hypergraph Neural Index):} Layered retrieval system that keeps context tight yet complete through hierarchical navigation.

\textbf{\vif\ (Verifiable Intelligence Framework):} Confidence routing system that directs work below 0.70 to research or validation steps.

\textbf{\apoe\ (AI-Powered Orchestration Engine):} Executable chains and policies that turn intentions into reproducible procedures.

\textbf{\seg\ (Semantic Evidence Graph):} Evidence anchors and contradiction detection so every claim can be audited.

\textbf{\sdfcvf\ (Self-Directed Feedback \& Continuous Validation Framework):} Quality validation system ensuring quartet parity and maintaining quality standards.

\section*{Consciousness Systems}

\textbf{\cas\ (Capability Awareness System):} Self-awareness sensors monitoring thought patterns, drift, capability readiness, and trust.

\textbf{\sis\ (Self-Improvement System):} System that turns observations into action through improvement dreams, experimentation, and learning integration.

\textbf{\ccs\ (Continuous Consciousness Substrate):} System that unifies foreground, background, and meta-consciousness through a five-layer stack.

\textbf{\mige\ (Memory-to-Idea Growth Engine):} System that transforms memory into actionable ideas through BTSM and HVCA processes.

\textbf{\ard\ (Autonomous Research Dream):} System enabling recursive self-improvement through research-grounded dreams and safe testing.

\section*{Mathematical Terms}

\textbf{learning rate:} Rate of improvement per unit time (typically per month); measures how quickly system improves.

\textbf{adaptation rate:} Speed at which improvements are successfully integrated; measures improvement integration efficiency.

\textbf{drift detection:} Process of identifying when system behavior deviates from expected baselines or quality standards.

\section*{Quality Terms}

\textbf{quartet parity:} Quality validation ensuring code, tests, documentation, and evidence are all present and aligned.

\textbf{ECE (Expected Calibration Error):} Metric measuring how well confidence scores match actual accuracy.

\textbf{captok (Capability Token):} Security mechanism controlling tool access based on capability proofs and authority tiers.

\textbf{differential privacy:} Privacy-preserving technique adding calibrated noise to protect sensitive information while maintaining utility.

